\section{把键盘输入送到用户态 Shell}

前几章分别完成了键盘中断、等待队列、文件描述符和控制终端。现在把这些机制
连起来。用户在 QEMU 窗口按下字母键时，i8042 产生扫描码和 IRQ1；内核解码
按键，把字符放入控制台输入队列，并唤醒正在等待读事件的前台线程。Shell
恢复运行后，通过 fd 0 读取这个字符。

\begin{pdfdiagram}
  \node[os hardware] (key) {QEMU 键盘\\PS/2 扫描码};
  \node[os block,right=of key] (irq) {i8042 + IRQ1\\Set 1 解码};
  \node[os state,right=of irq] (fifo) {控制台 FIFO\\256 字节};
  \node[os block,below=of fifo] (fd) {Shell 的 fd 0\\read/wait};
  \node[os state,left=of fd] (parser) {Ring 3 解析器\\固定容量};
  \node[os block,left=of parser] (service) {文件与目录\\系统调用};
  \draw[os arrow] (key) -- (irq);
  \draw[os arrow] (irq) -- (fifo);
  \draw[os arrow] (fifo) -- (fd);
  \draw[os arrow] (fd) -- (parser);
  \draw[os arrow] (parser) -- (service);
\end{pdfdiagram}
\diagramnote{宿主只产生虚拟按键；扫描码解码、排队、阻塞、唤醒和命令解析均由来宾软件完成。}

\topic{fd 0、1、2 让 Shell 不必认识设备端口}
Shell 启动时，fd 0 指向控制台输入，fd 1 和 fd 2 指向控制台输出。它只调用
描述符 ABI，不读取 \code{0x60}，也不向 \code{0x3F8} 写字符。驱动负责
设备寄存器，对象层负责权限和生命周期，用户程序只看到可以读写的小整数。

\begin{osgridlongtable}{L{0.11\textwidth}L{0.25\textwidth}L{0.20\textwidth}L{0.32\textwidth}}
\textbf{fd} & \textbf{当前对象} & \textbf{允许操作} & \textbf{Shell 的用途}\\ \hline
0 & ConsoleInput & read、wait readable & 读取命令行\\ \hline
1 & ConsoleOutput & write & 输出命令结果\\ \hline
2 & ConsoleError & write & 输出诊断信息\\ \hline
3 及以后 & 文件或目录对象 & 由打开标志决定 & \code{cat}、\code{write}、\code{ls}\\ \hline
\end{osgridlongtable}

控制台暂时没有数据时，读取包装先得到 would-block，再调用等待系统调用。
线程此时进入 Blocked；下一次输入到达后，它回到 Ready，重新尝试读取。
这里复用第 18 章已经讲过的 Try/Wait 循环，不再重复 WaitQueue 和 TTY
内部实现。

\topic{扫描码覆盖范围是可检查的实现边界}
当前键盘解码器处理 Shell 使用的 US 键区、Shift、Caps Lock、Enter 和
Backspace。Caps Lock 只影响字母；数字和标点只由 Shift 选择。Unicode、
输入法、布局切换和完整扩展键尚未实现，未识别的扫描码不会凭空生成字符。

控制台 FIFO 保存 256 字节，并分别统计成功提交、读取、丢弃和当前缓冲量。
在自动会话结束时，测试检查
\[
  \mathrm{submitted}
  =
  \mathrm{read}
  + \mathrm{buffered},
  \qquad
  0\le\mathrm{buffered}\le256.
\]
这个等式不能说明每个字符都正确，但能发现覆盖旧数据、重复读取和计数遗漏。

\section{最小运行时和 Shell 解析器怎样工作}

Shell 是普通 Ring 3 ELF。它没有 libc、iostream、异常、RTTI、堆分配或
全局动态构造。入口汇编准备参数后调用 C++ 函数；用户态系统调用包装按照项目
ABI 放置寄存器，再执行 \code{int 0x80}。

\topic{两套调用约定在汇编包装处相接}
普通 C++ 函数遵循 System V AMD64 ABI，前六个整数参数位于
RDI、RSI、RDX、RCX、R8、R9。项目系统调用 ABI 使用
RDI、RSI、RDX、R10 传递前四个参数。四参数包装把 R8 搬到 R10 后触发
软件中断；内核保存现场并按系统调用号分发。寄存器位置来自项目 ABI，不是
\code{int} 指令自动赋予的含义。

\begin{osgridlongtable}{L{0.22\textwidth}L{0.24\textwidth}L{0.36\textwidth}}
\textbf{边界} & \textbf{第四个参数} & \textbf{谁完成转换}\\ \hline
C++ 调用包装函数 & R8 & 编译器按照 System V ABI 生成调用\\ \hline
系统调用入口 & R10 & 用户汇编包装执行 \code{mov r10,r8}\\ \hline
内核分发器 & 保存帧中的 R10 & 内核按项目 ABI 读取\\ \hline
\end{osgridlongtable}

编译器还可能为聚合清零和复制生成 \code{memset}、\code{memcpy} 调用。
用户运行时显式实现这两个 C ABI 符号，并使用固定宽度类型遍历字节。ELF 审计
继续拒绝未解析符号和意外的 \code{.init\_array}；删除初始化节并不能代替
执行构造函数。

\topic{解析器把所有状态放在固定数组中}
行缓冲容量为 128 字节，最多保存八个参数。解析器维护读取位置、写入位置、
参数起点、引号状态和转义状态；每个参数只记录固定存储中的 offset 与 length。
这样一次失败不会留下引用旧行缓冲的悬空视图。

\begin{pdfdiagram}
  \node[os state] (normal) {普通状态};
  \node[os block,right=32mm of normal] (single) {单引号内};
  \node[os block,below=10mm of single] (double) {双引号内};
  \node[os hardware,below=10mm of normal] (escape) {下一字节转义};
  \draw[os arrow] ([yshift=2mm]normal.east) --
    node[above]{\scriptsize 单引号} ([yshift=2mm]single.west);
  \draw[os arrow] ([yshift=-2mm]single.west) --
    node[below]{\scriptsize 单引号} ([yshift=-2mm]normal.east);
  \draw[os arrow] (normal) -- node[left]{\scriptsize 反斜杠} (escape);
  \draw[os arrow] (normal) -- node[above,sloped]{\scriptsize 双引号} (double);
  \draw[os arrow] (double) -- node[below,sloped]{\scriptsize 双引号} (normal);
\end{pdfdiagram}
\diagramnote{单引号内反斜杠按普通字符处理；普通状态和双引号状态允许转义下一字节。}

解析器拒绝超长输入、参数过多、未闭合引号和末尾孤立反斜杠。输出对象在开始
解析时清零；错误返回时不会带出一条半完成命令。空行单独处理，不会落入
unknown command。

\topic{每条命令组合已经存在的系统调用}
\begin{osgridlongtable}{L{0.19\textwidth}L{0.34\textwidth}L{0.29\textwidth}}
\textbf{命令} & \textbf{主要调用} & \textbf{观察点}\\ \hline
\code{help}、\code{echo} & 写 fd 1 & 参数与有界输出\\ \hline
\code{pwd} & 写 fd 1 & 当前目录范围\\ \hline
\code{ls} & 打开目录、读取目录项、关闭 & 目录游标和 EOF\\ \hline
\code{mkdir} & 创建目录 & 路径复制和已存在错误\\ \hline
\code{write} & 打开、循环写入、关闭 & create、truncate、部分写\\ \hline
\code{cat} & 打开、循环读取、关闭 & 文件 EOF\\ \hline
\code{sync} & 同步文件系统 & 缓存写回和磁盘 FLUSH\\ \hline
\code{exit} & 退出进程 & fd 与地址空间回收\\ \hline
\end{osgridlongtable}

\code{pwd} 目前只输出根目录，因为系统还没有每进程当前工作目录和
\code{chdir}。这个限制由行为直接体现；Shell 不保存一份内核无法解析的
伪路径状态。

\section{用一次自动会话检查整条路径}

人工敲入命令适合初次观察，却不适合回归。QEMU runner 等待 Shell 输出
\code{READY}，再通过 QMP 的 \code{sendkey} 逐键发送下面的会话：

\begin{lstlisting}
help
echo hello world
pwd
mkdir /demo
write /demo/message hello
cat /demo/message
ls /demo
sync
unknown
exit
\end{lstlisting}

QMP 操作虚拟键盘前端。此后的扫描码、IRQ1、端口读取、字符队列、唤醒和
Ring 3 解析仍由来宾完成。测试没有调用内核的字符提交函数，也没有修改来宾
内存，所以硬件入口和等待路径仍在覆盖范围内。

\topic{检查结果，不固定合法的并发交错}
测试要求九条已知命令产生对应标记，unknown 被拒绝，Shell 最终退出。
\code{cat} 必须读到 \code{hello}，\code{ls} 必须读到
\code{message}。后台线程的串口输出可以在合法调度点交错，因此断言约束每个
执行流内部的顺序，不把某一次 QEMU 调度顺序写成 ABI。

\begin{osgridlongtable}{L{0.18\textwidth}L{0.30\textwidth}L{0.34\textwidth}}
\textbf{层级} & \textbf{输入} & \textbf{主要检查}\\ \hline
单元 & 扫描码、FIFO、解析字节 & 容量、状态转换和越界拒绝\\ \hline
集成 & 描述符、目录和用户 ELF & 权限、游标、ABI 与链接结构\\ \hline
随机 & 固定种子解析字节流 & 任意切片仍位于固定存储内\\ \hline
整机 & QMP 命令和同盘重启 & IRQ 到 Ring 3、持久化与回收\\ \hline
\end{osgridlongtable}

\topic{失败时先定位最后一个已出现的检查点}
完全没有 \code{READY} 时，先查看用户 ELF 是否装入、PID 1 是否进入
Running、标准 fd 是否建立。收到 \code{READY} 但没有命令字符时，再查看
IRQ1、扫描码解码和 FIFO 计数。命令开始执行却读不到文件时，检查路径解析、
文件描述符和文件系统日志。沿最后一个已经出现的日志向后找，比从 Shell
源文件盲目搜索更快。

\topic{当前用户环境仍保留明确边界}
系统已经能从磁盘装入程序、运行 Shell、使用管道和信号、控制前后台作业，并
通过 devfs 和 procfs 访问设备与系统状态。它仍缺少多用户权限、动态链接、
完整 POSIX 终端、Unicode 输入和 SMP。加入其中任何一项，都需要同时修改
用户 ABI、内核对象、失败路径和测试，不能只增加一个命令名。
